\section{Theoretical framework}
This research paper uses certain terms with ambiguous definitions. This section aims to explicitly state a definition of these terms.

\subsection{Partials, frequency ratios, chords, and harmonic complexity}
A [partial](FUNDAMENTAL) frequency is a fundamental value of a musical note, and holds the frequency which the note produces as its primary pitch. Over the course of this paper, a partial wave will be expressed both as a note on sheet music, and as a sine wave of the form $y = a\sin(2\pi * f* (x - d^\circ))$, where $a = \text{amplitude in }\unit{db}$, $f = \text{frequency in } \unit{hz}$, and $d = \text{ phase}$.

\begin{figure}
% figure of a C4 on sheet
% figure of a C4 as a wave
\end{figure}

A frequency ratio is a relation between two or more partial frequencies. An interval is a relation of two notes, while a triad is a relation of three notes. Over the course of this research paper, a frequency ratio is expressed as two or more notes on sheet music, and their corresponding distance on the musical scale (e.g. a "major third" or "perfect fifth"), or as a mathematical ratio of the form $a : b$ (for intervals) or $a : b : c$ (for triads), where $a$, $b$, and $c$ are partial frequencies, respective in size. 

\ref{The Arithmetic of Listening}

